% Introduction -- budget 682 words (body text).
% Sources: docs/paper_handoff.md sections 1, 4, 5, 7 (claims not to make); docs/results_summary.md (4.15-fold, 1.88 and
% 1.26 times the truth, 404 entries); Sections 2-7 of this paper. Research questions and hypotheses worded exactly as in
% Table 9 (D87). Every reference verified (notes in references.bib).
\section{Introduction}\label{sec:intro}

As distribution grids are digitalised, Morocco's among them \cite{elmaghraoui2026morocco}, feeder telemetry becomes the record from which outage time is counted and load is forecast. The public smart-meter records of feeders in four Moroccan cities \cite{bensalah2025uci,bensalah2025dib} have been used for load forecasting \cite{bensalah2025opsearch} and demand scheduling \cite{sarange2026}. Between the meter and the analyst, however, the data pass through a recording, export and processing chain that can fail: smart-meter outage notifications can go missing \cite{jiang2016outage}, and short gaps are commonly bridged by linear interpolation \cite{peppanen2016imputation}.

An anomaly in such a record is therefore ambiguous. A feeder that reads zero, or a series that stays flat or runs as a straight line, may show a real outage or a failure of the acquisition chain, and the two call for opposite treatment: an artefact should be imputed, a grid event kept. Counting every zero as an outage inflates outage statistics; interpolating every gap erases real outages. In the data studied here, counting every dead state as an outage gives 4.15 times the outage time of the states our method types as grid events; on injected data, where the truth is known, it gives 1.88 times the true outage time, and typed counting 1.26 times. Typing the anomalies decides whether outage statistics describe the grid or the data pipeline.

The records carry no event labels, a common situation in energy data \cite{fan2018autoencoder}. Multivariate detectors return an anomaly score, not a type \cite{hundman2018spacecraft,audibert2020usad,deng2021gdn}; data cleansing treats anomalies as errors to repair \cite{chen2010cleansing,mateos2013cleansing}; and bad data have been told from events in synchrophasor measurements of transmission grids \cite{wu2017lowquality,wu2021iforest}. Weak supervision has classified PMU events from a small labelled share once data quality was fixed \cite{liu2022weakly}; here, data quality is what must be labelled, and generic label models are unstable on the sparse votes that rules give (Section~\ref{sec:dsfail}).

This paper types each anomaly of a feeder record as an acquisition artefact or a grid event, with a calibrated probability, from the telemetry alone and without labelled data. The evidence is structural: the level of the hierarchy at which an anomaly appears (zone, city or region, with the acquisition chain as a fourth, non-electrical layer) and the traces that a recording chain leaves and a de-energised feeder does not. In the Moroccan records, by our typing, acquisition artefacts dominate: zero-fill gaps make up three quarters of the dead zone-hours, and the only regional-outage candidate is an acquisition failure. Straight-line interpolation is the most common fill, and the shape of a restoration carries almost no evidence: one dead run in five returns through a synthetic straight line. The contributions are:
\begin{itemize}
\item a typology of four artefact and five grid-event types, each with a measurable signature at one level of the hierarchy, including straight-line fills;
\item nineteen rules and a constrained Dawid--Skene label model that turns their sparse votes into calibrated probabilistic labels, stable across restarts where the unconstrained model and Snorkel's label model are not;
\item a multi-scale residual detector split into region, city and zone parts, which raises the recall of the rules at every contamination rate, at a cost in precision at low contamination;
\item an evaluation beyond detection scores: injected events that copy observed signatures, outage statistics and forecasts after cleaning, and transfer to a second system without re-tuning;
\item a typed catalogue of the dataset's anomalies (404 entries), released with the code, so that later studies can impute the artefacts and keep the grid events.
\end{itemize}

The paper asks five research questions:
\begin{itemize}
\item[RQ1] Can the hierarchy and the signatures alone, without labels, type the anomalies?
\item[RQ2] Which signatures separate grid events from artefacts?
\item[RQ3] Does learning add to the rules?
\item[RQ4] Does type-aware cleaning help downstream?
\item[RQ5] Does the method transfer without re-tuning?
\end{itemize}
It also tests three hypotheses on the physics of the signatures:
\begin{itemize}
\item[H1] Acquisition failures leave exact zeros, holds or straight lines; de-energised feeders leave a noise floor.
\item[H2] Outages restore with a ramp or an overshoot; gaps return directly.
\item[H3] Load transfers conserve current; faults and gaps do not.
\end{itemize}

Section~\ref{sec:related} reviews related work; Sections~\ref{sec:data}--\ref{sec:experiments} present the data and typology, the method and the experimental design. Section~\ref{sec:results} reports the results, Section~\ref{sec:discussion} answers the questions and states the limitations, and Section~\ref{sec:conclusions} concludes.
